\subsection{Rank Theorem}

\begin{theorem} \label{thm:constant-rank-theorem-rudin} \hlt{Rank Theorem}\\
Let $\mathbf{F}:E\subseteq\R^n\to\R^m$ be $C^1$ with constant rank $r$. At $\a\in E$, put $A=D\mathbf{F}(\a)$, $Y_1=\operatorname{range}(A)$, and let $P$ be the orthogonal projection onto $Y_1$ (Definition \ref{def:orthogonal-projection-subspace}). Locally, after a $C^1$ coordinate change $H$ with $C^1$ inverse on a small coordinate neighborhood $V_0$, there is a $C^1$ map $\varphi:A(V_0)\subseteq Y_1\to\ker(P)$ such that $\mathbf{F}(H(\x))=A\x+\varphi(A\x)$.\\
In coordinates adapted to $\ker A$ and $\operatorname{range}(A)$, this is the same local geometry as a projection onto $r$ coordinates.
\begin{proof}
If $r=0$, Theorem \ref{thm:bounded-derivative-convex-mean-value} makes $\mathbf{F}$ locally constant.\\
Assume $r>0$. Choose $S:Y_1\to\R^n$ with $AS=I_{Y_1}$, possible by choosing preimages of a basis of $Y_1$, and set $G(\x)=\x+SP(\mathbf{F}(\x)-A\x)$. Since $DG(\a)=I$, Theorem \ref{thm:inverse-function-theorem-rn} gives a local inverse $H=G^{-1}$.\\
Let $\Psi=\mathbf{F}\circ H$. From $AG=P\mathbf{F}$ we get $P\Psi(\x)=A\x$. Hence $PD\Psi(\x)=A$. The rank of $D\Psi(\x)$ is $r$ because $DH(\x)$ is invertible and $\mathbf{F}$ has constant rank $r$; the rank of $PD\Psi(\x)$ is also $r$.\\
Therefore $P$ is injective on $\operatorname{range}D\Psi(\x)$, and $D\Psi(\x)\mathbf{v}=0$ whenever $A\mathbf{v}=0$.\\
On a small convex coordinate neighborhood, any two points with the same $A\x$ differ by a vector in $\ker(A)$ and can be joined inside that affine fiber. The preceding derivative vanishing and Theorem \ref{thm:bounded-derivative-convex-mean-value} show that $\Psi$ is constant on such fibers. Thus $(I-P)\Psi(\x)$ depends only on $A\x$; call it $\varphi(A\x)$. Then $\mathbf{F}(H(\x))=\Psi(\x)=A\x+\varphi(A\x)$, and $\varphi$ is $C^1$ because $\Psi$ is $C^1$ and the dependence is through the coordinate $A\x$.\\
\hlt{Structural Commentary}: Constant rank gives a local graph normal form.\\
\hlt{Key Insight}: Same rank, same local level-set geometry.
\end{proof}
\end{theorem}

Rank-nullity (Theorem \ref{thm:rank-nullity}) says that the domain directions split into rank directions and fiber directions. The rank theorem says a constant-rank map has, locally, the same split after $C^1$ changes of coordinates.

\begin{figure}[H]
  \centering
  \begin{tikzpicture}[
  scale=0.9,
  >=Stealth,
  every node/.style={font=\scriptsize},
  fiber/.style={blue!70!black, line width=0.75pt},
  image/.style={black!70, line width=0.75pt},
  arrow/.style={->, line width=0.75pt},
  point/.style={circle, fill=black, inner sep=1.1pt}
]
  \begin{scope}
    \node[anchor=south] at (1.55,2.85) {original coordinates};
    \foreach \x/\s in {0.45/0.05,0.95/-0.04,1.45/0.08,1.95/-0.02,2.45/0.04} {
      \draw[fiber]
        (\x,0.20) .. controls (\x-0.16,0.85+\s) and (\x+0.18,1.35-\s) .. (\x,2.55);
    }
    \node[left=2pt] at (0.35,1.38) {fibers};
  \end{scope}

  \draw[arrow] (3.10,1.38) -- node[above=2pt] {$\mathbf{F}$} (4.18,1.38);

  \begin{scope}[xshift=4.45cm]
    \node[anchor=south] at (1.45,2.85) {image};
    \draw[image] (0.10,1.00) .. controls (0.90,1.50) and (1.85,1.35) .. (2.90,2.15);
    \draw[image] (0.20,1.30) .. controls (1.05,1.82) and (1.95,1.65) .. (3.02,2.48);
    \foreach \x/\y in {0.42/1.13,0.92/1.31,1.42/1.48,1.92/1.66,2.42/1.89} {
      \node[point] at (\x,\y) {};
    }
    \node[right=2pt] at (2.88,2.36) {$\operatorname{range}D\mathbf{F}$};
  \end{scope}

  \begin{scope}[yshift=-3.35cm]
    \node[anchor=south] at (1.55,2.85) {normal form};
    \foreach \x in {0.45,0.95,1.45,1.95,2.45} {
      \draw[fiber] (\x,0.20) -- (\x,2.55);
    }
    \node[left=2pt] at (0.35,1.38) {$\ker A$-fibers};
  \end{scope}

  \draw[arrow] (3.10,-1.97) -- node[above=2pt] {$P$} (4.18,-1.97);

  \begin{scope}[xshift=4.45cm,yshift=-3.35cm]
    \node[anchor=south] at (1.45,2.85) {projection image};
    \draw[image] (0.25,1.38) -- (3.00,1.38);
    \foreach \x in {0.45,0.95,1.45,1.95,2.45} {
      \node[point] at (\x,1.38) {};
    }
    \node[right=2pt] at (3.00,1.38) {$\R^r$};
  \end{scope}
\end{tikzpicture}

  \caption{After local coordinate changes, a constant-rank map looks like projection onto the rank directions.}
  \label{fig:rank-theorem-projection-normal-form}
\end{figure}

\begin{theorem} \label{thm:lagrange-multiplier-one-constraint} \hlt{One-Constraint Lagrange Multiplier}\\
Let $E\subseteq\R^n$ be open, $f,g:E\to\R$ be $C^1$, and $S=\{\x:g(\x)=c\}$.\\
If $p\in S$, $\nabla g(p)\neq\ze$, and $f|_S$ has a local extremum at $p$, then $\exists\lambda\in\R$ with $\nabla f(p)=\lambda\nabla g(p)$.
\begin{proof}
If $Df(p)=0$, take $\lambda=0$. Otherwise suppose $Df(p)$ is not a multiple of $Dg(p)$.\\
Complete $Dg(p),Df(p)$ to a dual basis $Dg(p),\quad Df(p),\quad \ell_3,\ldots,\ell_n$ by Theorem \ref{thm:dual-basis-and-dimension},\\
and set $\Phi(\x)=\bigl(g(\x)-c,\ f(\x)-f(p),\ \ell_3(\x-p),\ldots,\ell_n(\x-p)\bigr)$.\\
Then $D\Phi(p)$ is invertible, so Theorem \ref{thm:inverse-function-theorem-rn} makes $\Phi$ locally invertible. In these coordinates, $S$ is locally the slice with first coordinate $0$, while the second coordinate is $f-f(p)$. Local invertibility gives nearby points of $S$ with this second coordinate of both signs, contradicting extremality.\\
Hence $Df(p)=\lambda Dg(p)$; Riesz representation, Theorem \ref{thm:riesz-representation-finite-dimensional}, gives the gradient form.\\
\hlt{Structural Commentary}: The constraint kills first-order variation along $S$.\\
\hlt{Key Insight}: If $Df$ were independent of $Dg$, $f$ would be a free constraint coordinate.
\end{proof}
\end{theorem}

\begin{method} \label{meth:several-variable-differentiation-proof-patterns-source-todo}\label{meth:several-variable-differentiation-proof-patterns} \hlt{Several-Variable Differentiation Proof Patterns}
\begin{enumerate}[label=(\roman*)]
\setlength{\itemsep}{0pt}
\item \hlt{Linear approximation}: write $\f(\x+\h)-\f(\x)=D\f(\x)\h+\mathbf{r}(\h)$ with $\norm{\mathbf{r}(\h)}_2=o(\norm{\h}_2)$.
\item \hlt{Norm estimates}: use $\norm{A\h}_2\leq\normop{A}\norm{\h}_2$ and the product estimate from Proposition \ref{prop:spectral-norm-basic-properties}.
\item \hlt{Line-segment reduction}: on convex sets, restrict to $\gamma(t)=(1-t)\a+t\b$ and apply one-variable results.
\item \hlt{Local inversion}: if $A=D\f(\a)$ is invertible, write $D\f(\x)=A+E_\x$ with $\normop{E_\x}$ small and solve by contraction (Theorem \ref{thm:contraction-principle}); linear block equations use Theorem \ref{thm:linear-implicit-function-theorem}.
\item \hlt{Coordinates}: read partial derivatives as the entries of the matrix of $D\f(\x)$, using Theorem \ref{thm:total-derivative-matrix-partials}.
\item \hlt{Higher derivatives}: regard $D^r\f(p)$ as an $r$-linear map; recover coordinate partials via Theorem \ref{thm:second-derivative-partials-and-continuity}.
\item \hlt{Rank normal forms}: for $A=D\mathbf{F}(\a)$, split domain by $\ker A$ and image by $\operatorname{range}A$; use the quotient-by-kernel picture from Theorem \ref{thm:quotient-by-kernel-is-range} and project with $P_{\operatorname{range}A}$.
\end{enumerate}
\hlt{Key Insight}: Several-variable calculus is one-variable limiting plus finite-dimensional linear algebra.
\end{method}
